\documentclass[10pt]{article}

\usepackage[utf8]{inputenc}

\usepackage[T1]{fontenc}

\usepackage{amsmath}

\usepackage{amsfonts}

\usepackage{amssymb}

\usepackage[version=4]{mhchem}

\usepackage{stmaryrd}

\usepackage{bbold}

\usepackage{graphicx}

\usepackage[export]{adjustbox}

\graphicspath{ {./images/} }



\begin{document}

ным алгебраич. многообразием. Подгруппа $P \subset G$ тогда и только тогда является П. І., когда она содержит какую-нибудь Бореля подгрупny группы $G$. $\Pi$ а р аболической шод г ру у ІІ пй группы $G_{k} k$ рациональных точек группы $G$ наз. подгруппа $P_{k} \subset G_{k}$, являющаяся группой $k$-рациональных точек нек-рой П. П. $\boldsymbol{P}$ в $G$ и плотная в $\boldsymbol{P}$ в тошологии Зариского. Если $\operatorname{char} k=0$ и $\mathfrak{g}$ - алгебра Ли группы $G$, то замкнутая подгруппа $P \subset G$ тогда и только тогда является П. П., когда ее алгебра Ли является параболической подалгеброй $\boldsymbol{8} \mathfrak{g}$.



Пусть $G$ - связная редуктивная линейная алгебраич. группа. Минимальные $k$-замквутые (т. е. определенные над $k$ ) П. П. играют в теории таких групп для произвольного основного поля $k$ ту же роль, что подгруппы Бореля для алгебраически замкнутого поля $k$ (см. [1]). В частности, любые две минимальные $k$-замкнутые П. п. группы $G$ сопряжены над $k$. Если две $k$-замкнутые П. п. группы $G$ сопряжены над каким-нибудь расширением поля $k$, то они сопряжены и над $k$. Множество классов сопряженных П. П. (соответственно множество классов сопряженных $k$-замкнутых П. п.) группы $G$ состоит из $2^{r}$ (соответственно $2^{r_{k}}$ ) элементов, где $r-$ ранг коммутанта $(G, G)$ группы $G$, а $r_{k}$ - его $k$-ранг, равный размерности максимального расщепимого над $k$ тора в $(G, G)$. Точнее, каждый такой класс определяется нек-рым произвольным подмножеством множества простых корней (соответственно простых $k$-корней) группы $G$ аналогично тому, как каждая параболич. подалгебра редуктивной алгебры Ли сопряжена одной из стандартных подалгебр (см. [2], [4]).



Всякая П. п. $P$ группы $G$ связна, совпадает со своим нормализатором и допускает р а 3 л о ж е н и е Л е в и, т. е. может быть представлена в виде полуІрямого произведения своего унипотентного радикала и нек-рой $k$-замкнутой редуктивной подгруппы, наз. п о д г р у п п о й Л е в и группы $P$. Любые две подгруппы Леви в П. П. $P$ сопряжены посредством рационального над $k$ элемента группы $P$. Две П. П. группы $G$ наз. п ро т и в о п о о ж ны ми, если их пересечение является подгруппой Леви в каждой из них. Замкнутая подгруппа группы $G$ тогда и только тогда является П. П., когда она совпадает с нормализатором своего унипотентного радикала. Всякая максимальная замкнутая подгруппа группы $G$ либо является П. п., либо имеет редуктивную связную компоненту единицы (cM. [2], [4]).



П. П. в группе $\mathrm{GL}_{n}(k)$ невырожденных линейных преобразований $n$-мерного векторного пространства $V$ над полем $k$ исчерпываются подгруппами $P(v)$, состоящими из всех автоморфизмов пространства $V$, которые сохраняют фиксированный флаг типа $v=\left(n_{1}, n_{2}, \ldots, n_{t}\right)$ в $V$. При этом факторпространство GI. $(k) / P(v)$ совпадает с многообразием всех фллагов типа $v$ в пространстве $V$.



В случае, когда $k=\mathbb{R}, \mathbb{R}$-замкнутые П. п. допускают следующую геометрич. интерпретацию (см. [5]). Пусть $G_{\mathbb{R}}$ - некомпактная полупростая вещественная группа Ли, совпадаюпая с группой вещественных точек полупростой алгебраич. группы $G$. Подгруппа группы $G_{\mathbb{R}}$ тогда и только тогда является П. П., когда она совпадает с группой движений соответствующего некомпактного симметрич. пространства $M$, сохраняющих нек-рый $k$-пучок геодезич. лучей в $M$ (два геодезич. луча в $M$ считаются принадлежащими одному $k$-пучку, если расстояние между двумя точками, движущимися с одинаковыми постоянными скоростями вдоль этих лучей в бесконечность, стремится к конечному пределу).



\begin{enumerate}

  \setcounter{enumi}{1}

  \item П. п. с и с т е м ы Т и т с а $(G, B, N, S)-$ подгруппа группы $G$, сопряженная подгруппе, содержащей $B$. Каждая П. П. совпадает со своим нормализатором. Пересечение любых двух П. П. содержит подгруппу группы $G$, сопряженную с $T=B \cap N$. В частности, П. п. системы Титса, связанной с редуктивной линейной алгебраич. грушпой $G,-$ это то же, что П. П. группы $G$ (см. [3], [4]).

\end{enumerate}



Лит.: [1] Б о р е л ь А., Т и т с Ж., “Математика», 1967, т. $11, \mathcal{N}_{2} 1$, c. $43-111$; [2] и $\times$ же, там же, 1972, т. 16, Nㅡ 3 , c. $3-12 ;[3]$ Б у р б а к и Н., Группы и алгебры Ли. Подалгебры Картана, регулярные әлементы, расщепляемые полупростые алгебры Ли, пер. с франц., М., 1978; [4] Х а м ф р и Д Ж. Линейные алгебраические группы, пер. с англ., М., 1980; ${ }^{2}[5]$ К а p п е л е в и ч Ф. И., "Труды Моск. матем. об-ва", 1965

В. Л. Попов



ПАРАБОЛИЧЕСКАЯ РЕГРЕССИЯ, П $о$ Л и н $о$ м иа ль в а я е г р с с и я,- модель регрессии, в к-рой функции регрессии суть многочлены. Точнее, пусть $\quad X=\left(X_{1}, \ldots, X_{m}\right)^{T}$ и $Y=\left(Y_{1}, \ldots, Y_{n}\right)^{T}-$ случайные векторы, принимающие значения $x=\left(x_{1}, \ldots\right.$, $\left.x_{m}\right)^{T} \in \mathbb{R}^{\boldsymbol{m}}$ и $y=\left(y_{1}, \ldots, y_{n}\right)^{\boldsymbol{T}} \in \mathbb{R}^{n}$, и пусть существует\[

E\{Y \mid X\}=f(X)=\left(f_{1}(X), \ldots, f_{n}(X)\right)^{T}

\]



(т. е. супествуют $E\left\{Y_{1} \mid X\right\}=f_{1}(X), \ldots, E\left\{Y_{n} \mid X\right\}=$ $\left.=f_{n}(X)\right)$. Регрессия наз. п а р а б о ли ч е с к о й, если компоненты вектора $\mathrm{E}\{Y \mid X\}=f(X)$ суть многочлены от компонент вектора $X$. Напр., в простейшем случае, когда $Y$ и $X$ - обычные случайные величины, уравнение П. р. имеет вид



\[

y=\beta_{0}+\beta_{1} X+\ldots+\beta_{p} X^{p}

\]



где $\beta_{0}, \ldots, \beta_{p}-$ коэфффициенты регрессии. Частный случай П. р.- линейная регрессия. Добавлением к вектору $X$ новых компонент можно всегда свести П. p. к линейной. См. Регрессия, Регрессионный анализ.



Лит.: [1] К р а м е р Г., Математические методы статистики, пер. с англ., 2 изд., М., 1975; [2] С е б е р Д ж., Линейный регрессионный анализ, пер. с англ., М., 1980.



ПАРАБОЛИЧЕСКАЯ СПИРАЛЬ - плоская Транс цендентная кривая, уравнение к-рой в полярных координатах имеет вид



\[

\rho=a \sqrt{\bar{\varphi}}+l, \quad l>0

\]



Каждому значению $\varphi$ соответствуют два значения $\sqrt{\varphi}$ - положительное и отрицательное. Кривая имеет бесконечно много двойных точек и одну точку перегиба (см.



\includegraphics[max width=\textwidth, center]{2023_07_08_e97951c69e21fdb964e9g-1(1)}

рис.). Если $l=0$, то кривая наз. Ферма спиралью. П. с. относится к т. н. алгебраич. спиралям (см. $\mathrm{Cnu}-$ рали). Лum.: [1] С а в е ло в А. А., Плоские кривые, М., 1960.



ПАРАБОЛИЧЕСКАЯ ТОЧКА - точка регулярной поверхности, в к-рой соприкасающийся параболоид вырождается в параболич. цилиндр. В П. т. Дюпена индикатриса является парой параллельных прямых, гауссова кривизна равна нулю, одна из главных кривизн обрацается в нуль, а для коәфффициентов второй квадратичной формы справедливо равенство



\[

L N-M^{2}=0 .

\]



ПАРАБОЛИЧЕСКИЕ КООРДИНАТЫ - числа $u$ и $v$, связанные с прямоугольными координатами $x$ и $y$ формулами



\[

x=u^{2}-v^{2}, \quad y=2 u v,

\]



где $-\infty<u<\infty, 0 \leqslant v<\infty$. Координатные линии: две системы взаимно ортогональных парабол с противоположно направленными осями.



Коэфффициенты Ламе:



\[

L_{u}=L_{v}=2 \sqrt{u^{2}+v^{2}} .

\]



Д. Д. Соколов.



\begin{center}

\includegraphics[max width=\textwidth]{2023_07_08_e97951c69e21fdb964e9g-1}

\end{center}





\end{document}